\chapter{Linear ODEs}

\section{A single linear equation}
\[
\dot x = ax, \qquad x(0)=x_0.
\]
Solve: $\frac{d}{dt}\log x = \dot x/x = a$, so $\log x = \log x_0 + at$, i.e.
\[
x(t) = x_0 e^{at}.
\]
$a>0$: exponential growth. $a<0$: exponential decay. $a=0$: constant.

\section{Why linearity matters}
A system $\dot x = Ax$ ($x\in\mathbb{R}^n$, $A$ a matrix) is
\emph{linear} in $x$. Its key property: if $x_1(t), x_2(t)$ both solve
$\dot x = Ax$, so does any combination $c_1x_1+c_2x_2$, since
\[
\frac{d}{dt}(c_1x_1+c_2x_2) = c_1\dot x_1 + c_2\dot x_2 = c_1Ax_1+c_2Ax_2 = A(c_1x_1+c_2x_2).
\]

\section{Diagonal systems in two dimensions}
A \emph{diagonal} system decouples into two copies of
Section~1:
\[
\dot x = \lambda_1 x, \qquad \dot y = \lambda_2 y
\quad\Longrightarrow\quad
x(t)=x_0e^{\lambda_1 t},\ \ y(t) = y_0e^{\lambda_2 t}.
\]
The qualitative picture depends on the signs of $\lambda_1,\lambda_2$:

\begin{figure}[h]
\centering
\includegraphics[width=\linewidth]{visualizations/phase_portraits_real.pdf}
\caption{Both negative: every trajectory decays to the origin
  (\emph{stable node}). Both positive: every trajectory grows away
  from it (\emph{unstable node}). Opposite signs: the origin attracts
  along one axis and repels along the other (\emph{saddle}).}
\end{figure}

\section{Diagonal systems with complex eigenvalues}
Take $A = \begin{pmatrix}a&-b\\b&a\end{pmatrix}$, i.e.\
$\dot x = ax-by,\ \dot y = bx+ay$. Its eigenvalues solve
$\det(A-\lambda I)=(a-\lambda)^2+b^2=0$, so $\lambda = a\pm ib$:
complex. Set $z=x+iy$; then
\[
\dot z = \dot x+i\dot y = (ax-by)+i(bx+ay) = (a+ib)(x+iy) = (a+ib)z,
\]
the single linear equation of Section~1, now with complex $a+ib$.
Its solution is $z(t) = z_0e^{(a+ib)t} = z_0e^{at}\big(\cos bt +
i\sin bt\big)$: a rotation at rate $b$, scaled by $e^{at}$.

\begin{figure}[h]
\centering
\includegraphics[width=\linewidth]{visualizations/phase_portraits_complex.pdf}
\caption{$a<0$: spiral sink. $a=0$: pure rotation, a \emph{center}.
  $a>0$: spiral source.}
\end{figure}

\section{Every linear system is diagonal, in the right coordinates}
Let $A$ have eigenvectors $v_1,v_2$ (eigenvalues $\lambda_1,\lambda_2$,
possibly complex conjugates) and $P=\begin{pmatrix}v_1&v_2\end{pmatrix}$
as in the eigendecomposition $A=PDP^{-1}$. Write $x = Pu$, i.e.\ $u =
P^{-1}x$ gives the coordinates of $x$ in the eigenbasis. Then
\[
\dot u = P^{-1}\dot x = P^{-1}Ax = P^{-1}APu = Du.
\]
So $u$ solves a \emph{diagonal} system --- Section~3 if
$\lambda_1,\lambda_2$ are real, Section~4 if they are complex
conjugates --- and $x(t) = Pu(t)$ translates the solution back to the
original coordinates. Every linear system is, up to a change of basis,
one of the pictures above.
